\section{Local endpoint maps and their amplification}\label{sec:maps}

In this section we prove \cref{prop:amplification}. We assume \cref{ass:G,ass:U}. Let
$\varepsilon\in(0,1/2)$ be the parameter of \cref{def:good}, to which all caps and $M$-goodness refer
(\cref{def:good,def:map}), and assume that $\EE_{p_c}\abs{C_o}^\alpha=\infty$ for some
$\alpha\in(0,1/2-\varepsilon)$. We fix $\beta\in(\alpha,1/2-\varepsilon)$, so that
\begin{equation}\label{eq:beta-sigma}
  \beta+\sigma<1\,,\qquad\sigma=\tfrac12+\varepsilon\,.
\end{equation}
Configurations have parameter $p_c$ unless stated otherwise.

The heavy tail provides infinitely many scales $n$ at which clusters of size about $n$ are frequent
(\cref{sec:scales}). One amplification step takes a map $\Phi$ and such a scale~$n$ and constructs a new
map, in four parts: a splicing experiment (\cref{sec:splicing}), an inverse description of its outcome
(\cref{sec:codes}), first and second moment bounds for the resulting density (\cref{sec:load}), and a
truncation (\cref{sec:iteration}). Iterating the step proves \cref{prop:amplification}.

\subsection{Heavy scales and a first map}\label{sec:scales}

\begin{lemma}[Heavy scales]\label{lem:scales}
There is an infinite set $\mathcal N$ of powers of two such that for every $n\in\mathcal N$
\begin{equation}\label{eq:scales}
  \pi_n\coloneqq\PP_{p_c}(n\le\abs{C_o}<2n)\ge n^{-\beta}
  \qquad\text{and}\qquad
  \PP_{p_c}(n/2\le\abs{C_o}<2n)\le(1+2^\beta)\,\pi_n\,.
\end{equation}
\end{lemma}

\begin{proof}
Put $\pi'_j\coloneqq\PP_{p_c}(2^j\le\abs{C_o}<2^{j+1})$ for $j\ge0$. Since $\abs{C_o}$ is almost surely
finite (\cref{lem:pc}), $\EE_{p_c}\abs{C_o}^\alpha\le\sum_j2^{\alpha(j+1)}\pi'_j$. If the sequence
$2^{\beta j}\pi'_j$ were bounded, the right-hand side would be finite because $\alpha<\beta$, contrary to
our assumption. So the sequence $2^{\beta j}\pi'_j$ is unbounded, and it has infinitely many
record indices~$j$, that is, $2^{\beta j}\pi'_j>2^{\beta i}\pi'_i$ for all $i<j$. Keep the records with
value at least one. At such a record, $\pi'_j\ge2^{-\beta j}$ and $\pi'_{j-1}<2^\beta\pi'_j$, which is
\eqref{eq:scales} for $n=2^j$.
\end{proof}

The next lemma collects the basic properties of local endpoint maps. For a map $\Phi$ and $p\in[0,1]$
we call $F_p(x,y)\coloneqq\EE_p\Phi_x(\omega,y)$ its \emph{unconditional endpoint kernel} at
parameter~$p$, and $s_p\coloneqq\sum_yF_p(x,y)$ its success probability at~$p$; thus $s=s_{p_c}$ and
$K=F_{p_c}/s$.

\begin{lemma}[Basic properties]\label{lem:map-basic}
Let $\Phi$ be a local endpoint map with $s>0$. Then $K$ is invariant and doubly stochastic, and
$r=\norm K\le1$. For every $p\in[0,1]$ the kernel $F_p$ is invariant, $F_p(x,y)=0$ if
$\dist(x,y)>r_\Phi+1$, its row and column sums equal $s_p$, and $\norm{F_p-F_{p_c}}\to0$ as
$p\to p_c$.
\end{lemma}

\begin{proof}
Invariance of $F_p$ follows from the covariance of $\Phi$ and the invariance of $\PP_p$. If
$\dist(x,y)>r_\Phi+1$, closing all edges without an endpoint in $B(x,r_\Phi)$ disconnects $x$ from~$y$
without changing $\Phi_x(\omega,\cdot)$, so $\Phi_x(\omega,y)=0$ for all~$\omega$. By the mass-transport
principle the column sums of the nonnegative invariant kernels $F_p$ and $\abs{F_p-F_{p_c}}$ equal their
row sums; in particular $K$ is doubly stochastic, and the Schur test gives $r\le1$. The entries of $F_p$
are polynomials in $p$, since $\Phi_x$ depends on finitely many edges. Hence the row sums of
$\abs{F_p-F_{p_c}}$, which have at most $\abs{B(o,r_\Phi+1)}$ nonzero terms, tend to zero, and the
Schur test gives $\norm{F_p-F_{p_c}}\to0$.
\end{proof}

\begin{lemma}[A first map]\label{lem:start}
For every $M_1$ there is a local endpoint map with cap $M\ge M_1$ and $r<\e^{-1}$.
\end{lemma}

\begin{proof}
Fix $\ell_0\in\N$ with $2\rho^{\ell_0}<\e^{-1}$. For a path $w$ of $\ell_0$ steps of the lazy walk, let $e(w)$ be
the number of distinct edges it traverses; holding steps traverse no edge. Call $w$ \emph{open} if all
edges it traverses are open. For $M\ge M_1$ put
\begin{equation*}
  \Phi_x(\omega,y)\coloneqq\1\{\abs{C_x}\le M,\ C_x\text{ is $M$-good}\}
  \sum_{w\colon x\to y}\mu(w)\,p_c^{\ell_0-e(w)}\,\1\{w\text{ is open}\}\,,
\end{equation*}
where the sum runs over $\ell_0$-step paths from $x$ to~$y$. This is a local endpoint map with cap $M$:
the sum over $y$ is at most $\sum_w\mu(w)=1$, an open path joins $x$ to~$y$, and by \cref{prop:good} the
indicator depends only on the edges near~$x$. Without the indicator, the unconditional endpoint
kernel would be $\sum_w\mu(w)p_c^{\ell_0-e(w)}p_c^{e(w)}=p_c^{\ell_0} P^{\ell_0}(x,y)$. Hence $F_{p_c}\le
p_c^{\ell_0} P^{\ell_0}$ entrywise and
\begin{equation*}
  s\ge p_c^{\ell_0}-\PP_{p_c}(\abs{C_o}>M)-\PP_{p_c}(\abs{C_o}\le M,\ C_o\text{ not $M$-good})\ge p_c^{\ell_0}/2
\end{equation*}
for $M$ large, by \cref{lem:pc,prop:good}. Thus $r=\norm{F_{p_c}}/s\le p_c^{\ell_0}\rho^{\ell_0}/(p_c^{\ell_0}/2)
=2\rho^{\ell_0}<\e^{-1}$.
\end{proof}

\subsection{The splicing experiment}\label{sec:splicing}

For the rest of this section fix an old map $\Phi$ with cap $M\ge2$, success probability~$s$, endpoint
kernel $K$ and $r=\norm{K}$; fix $\ell\in\N$ and $n\in\mathcal N$ with $\ell M\le n/2$, and put
$\nu\coloneqq \ell M$. For vertices $x$ and $z$ we run the following experiment.
\begin{enumerate}[label=(\arabic*)]
  \item Sample independent configurations $\omega^{(0)},\omega^{(1)},\dots,\omega^{(\ell)}$ with law
  $\PP_{p_c}$. We call $\omega^{(0)}$ the \emph{distinguished} configuration.
  \item Starting from $y_\ell\coloneqq z$, for $i=\ell,\ell-1,\dots,1$ apply the rule $\Phi_{y_i}$ to
  $\omega^{(i)}$, using randomness that is independent across the steps and of all other randomness;
  if it outputs a vertex, call it $y_{i-1}$, and if it fails, the experiment fails.
  \item The experiment \emph{succeeds} if all $\ell$ rules succeed, the cluster
  $A\coloneqq C_x(\omega^{(0)})$ satisfies $n\le\abs{A}<2n$, and $y_0\in A$.
\end{enumerate}
A successful experiment has \emph{weight}~$1/n$. Let $C_i\coloneqq C_{y_i}(\omega^{(i)})$ be the cluster
used in the $i$-th step; it contains $y_{i-1}$, has at most $M$ vertices and is $M$-good.

On success, the following \emph{switching procedure} splices pieces of the clusters $C_i$ onto $A$. It
maintains a set $S$, initially $S=A$, and a set of eligible indices, initially $\{1,\dots,\ell\}$. While
$z\notin S$, it performs a round:
\begin{enumerate}[label=(\roman*)]
  \item let $i$ be the largest eligible index with $C_i\cap S\neq\emptyset$, and declare all indices
  $\le i$ ineligible;
  \item let $B_i$ be the component of $y_i$ in the subgraph of $\omega^{(i)}$-open edges induced by
  $V\setminus S$, and let $\mathcal K_i$ be the set of $\omega^{(i)}$-open edges between $B_i$ and~$S$;
  \item choose the \emph{seam} $e_i\in\mathcal K_i$ uniformly at random, independently of everything
  else;
  \item exchange the states of $\omega^{(0)}$ and $\omega^{(i)}$ on the \emph{mask} $\mathcal M_i$, which
  consists of $e_i$ and of all edges that have an endpoint in $B_i$ and no endpoint in~$S$;
  \item replace $S$ by $S\cup B_i$.
\end{enumerate}
The \emph{output} consists of the configurations $\omega'^{(0)},\dots,\omega'^{(\ell)}$ after all
exchanges. The random objects of a successful run, namely the input configurations, the chain
$y_0,\dots,y_{\ell-1}$ and the seams, form its \emph{trajectory}~$\tau$.

\begin{lemma}[Geometry of the switching procedure]\label{lem:switching}
For a successful experiment every round is well defined, and the procedure stops after at most $\ell$
rounds. Let $i_1<\dots<i_t$ be the selected indices. Before each round $S=C_x(\omega^{(0)})$ for the
current $\omega^{(0)}$, and after the round $C_x(\omega^{(0)})=S\cup B_i$, with $e_i$ the only open edge
between $S$ and~$B_i$. The piece $B_{i_1}$ is attached to $A$ by $e_{i_1}$, and each later piece
$B_{i_j}$ is attached to $B_{i_{j-1}}$ by $e_{i_j}$ and to no other earlier piece. Consequently, in the
final distinguished configuration $\omega'^{(0)}$,
\begin{equation}\label{eq:switch-cluster}
  C_x(\omega'^{(0)})=A\cup B_{i_1}\cup\dots\cup B_{i_t}\ni z\,,\qquad
  \abs{C_x(\omega'^{(0)})}<2n+\nu\,,
\end{equation}
the seams are bridges of $C_x(\omega'^{(0)})$ separating $x$ from $z$, in the order
$e_{i_1},\dots,e_{i_t}$ from~$x$, and the open edges of $\omega'^{(0)}$ inside each $B_i$ are those
of~$\omega^{(i)}$. Every edge whose state in the distinguished configuration is changed has an endpoint
in the \emph{hanging part} $H\coloneqq B_{i_1}\cup\dots\cup B_{i_t}$, which is the vertex set of the
component of $C_x(\omega'^{(0)})$ beyond the first seam $e_{i_1}$.
\end{lemma}

\begin{proof}
Initially $C_1\ni y_0\in A$, so an index meets~$S$. Suppose index $i$ is selected. Then $y_i\notin S$:
for $i<\ell$ this holds because $C_{i+1}\ni y_i$ and $i+1$ was eligible but not selected, and for $i=\ell$
because $y_\ell=z\notin S$. Since $C_i$ is connected, meets $S$ and contains~$y_i$, the piece
$B_i\subseteq C_i\setminus S$ has at least one open edge to~$S$, so $\mathcal K_i\neq\emptyset$.

Before the exchange all edges of the current $\omega^{(0)}$ leaving $S$ are closed. The exchange
imports the $\omega^{(i)}$-states of all edges inside $B_i$ and of all edges from $B_i$ to
$V\setminus(S\cup B_i)$; the latter are closed in~$\omega^{(i)}$, because $B_i$ is a component
of the $\omega^{(i)}$-open subgraph on $V\setminus S$. It also opens~$e_i$. All other edges between
$B_i$ and $S$ keep their closed states, and no edge inside $S$ changes. Hence the new cluster of $x$ is
$S\cup B_i$, with $e_i$ its only open edge between $S$ and~$B_i$. Edges inside $B_i$ are not touched
again, since later masks consist of edges with no endpoint in the enlarged~$S$, apart from their seams.

Every eligible index larger than the selected $i$ missed the old~$S$. Hence the next selected cluster
meets the enlarged set only in~$B_i$, so the next piece is attached to $B_i$ and has no open edge to
earlier pieces or to~$A$. If $i<\ell$ and $z\notin S\cup B_i$, then $C_{i+1}\ni y_i\in B_i$ provides an
eligible index for the next round; if $i=\ell$, then $z=y_\ell\in B_\ell$ and the procedure stops. Since the
selected indices increase, there are at most $\ell$ rounds. The pieces form a chain attached by single
seams, which gives the claims about the seams and~\eqref{eq:switch-cluster}, using $\abs A<2n$ and
$\abs{B_i}\le M$. Every mask edge has an endpoint in its piece, which proves the last claim.
\end{proof}

\begin{figure}[t]
\centering
% The switched root cluster as a chain of pieces attached by single seams (schematic)
\begin{tikzpicture}[scale=0.95,
    vtx/.style={circle, fill, inner sep=1.3pt},
    piece/.style={draw=figaux, fill=figaux!10, rounded corners=9pt, thick},
    root/.style={draw=figgrey, fill=figgrey!12, rounded corners=12pt},
    seam/.style={figmain, line width=1.6pt}]
  \draw[root] (0,-0.9) rectangle (3.2,0.9);
  \node[font=\small] at (1.6,-1.25) {$A$};
  \node[vtx] (x) at (0.6,0) {};
  \node[font=\small, anchor=east] at (x.west) {$x$};
  % pieces
  \draw[piece] (4.2,-0.6) rectangle (6.2,0.6);
  \draw[piece] (7.2,-0.6) rectangle (9.2,0.6);
  \draw[piece] (11.0,-0.6) rectangle (13.0,0.6);
  \node[font=\small] at (5.2,-0.95) {$B_{i_1}$};
  \node[font=\small] at (8.2,-0.95) {$B_{i_2}$};
  \node[font=\small] at (12.0,-0.95) {$B_{i_t}$};
  \node at (10.1,0) {$\cdots$};
  % seams
  \draw[seam] (3.0,0.2) -- (4.4,0.1);
  \draw[seam] (6.0,-0.1) -- (7.4,0.0);
  \draw[seam] (9.0,0.15) -- (9.55,0.15);
  \draw[seam] (10.65,0.1) -- (11.2,0.1);
  \node[font=\small, text=figmain, anchor=south] at (3.7,0.2) {$e_{i_1}$};
  \node[font=\small, text=figmain, anchor=south] at (6.7,0.05) {$e_{i_2}$};
  \node[font=\small, text=figmain, anchor=south] at (10.9,0.15) {$e_{i_t}$};
  \node[vtx] (z) at (12.6,0) {};
  \node[font=\small, anchor=west] at (z.east) {$z$};
\end{tikzpicture}

\caption{The cluster of $x$ after the switching procedure (schematic). Each inserted piece $B_{i_j}$ is
attached to its predecessor by a single open seam $e_{i_j}$, so the seams are bridges separating $x$
from~$z$. Cutting the seams recovers $A$ and all pieces (\cref{lem:switching,lem:codes}).}\label{fig:chain}
\end{figure}

\subsection{Undoing the splicing}\label{sec:codes}

Fix $x$ and~$z$. For a configuration $\omega'$ with $z\in C_x(\omega')$, list the bridges of
$C_x(\omega')$ that separate $x$ from $z$ in their order from $z$ towards~$x$. A \emph{code} is a
sequence $c=(c_\ell,\dots,c_1)$ with entries in $\{\ast\}\cup\{1,\dots,M\}$. A successful run with output
distinguished configuration $\omega'^{(0)}$ has the following code: $c_i=\ast$ if $i$ is not selected,
and for the selected indices, read from $z$ towards $x$, the entry $c_{i_j}$ is the number of
separating bridges after the previous seam (or after $z$), up to and including $e_{i_j}$. We call
$c_{i_j}$ the \emph{gap} of $e_{i_j}$. Recall that $\sigma=1/2+\varepsilon$, and put
\begin{equation*}
  \Theta_M\coloneqq1+\sum_{m=1}^M\min\Bigl\{1,\frac{2M^\sigma}{m}\Bigr\}\le CM^\sigma\log M\,.
\end{equation*}

\begin{lemma}[Inverse codes]\label{lem:codes}
Fix $x$, $z$ and output configurations $\omega'^{(0)},\dots,\omega'^{(\ell)}$.
\begin{enumerate}[label=\textup{(\alph*)}]
  \item\label{it:codes-inverse} A code determines the input configurations, the sets $A$ and $B_i$, the
  seams, the masks and the contact sets $\mathcal K_i$ of every successful run with this output and
  this code.
  \item\label{it:codes-gaps} In every successful run each gap satisfies
  $1/\abs{\mathcal K_{i}}\le\min\{1,2M^\sigma/c_{i}\}$.
  \item\label{it:codes-sum} Summed over the codes $c$ of successful runs with this output,
  $\sum_c\prod_{i\colon c_i\neq\ast}\abs{\mathcal K_i}^{-1}\le\Theta_M^\ell$.
\end{enumerate}
\end{lemma}

\begin{proof}
\ref{it:codes-inverse}: The code locates the seams among the separating bridges. By
\cref{lem:switching}, cutting them decomposes $C_x(\omega'^{(0)})$ into $A$ (the part containing $x$)
and the pieces $B_{i_1},\dots,B_{i_t}$ in their order. This determines the set $S$ before each round,
hence every mask, and undoing the exchanges in reverse order recovers the input configurations. The
contact sets are then determined by the input configurations and the pieces.

\ref{it:codes-gaps}: The separating bridges counted in the gap of $e_i$, other than $e_i$ itself, lie
inside $B_i$, between the point where the next seam (or $z$) attaches and the endpoint of $e_i$.
Since $B_i$ is joined to the rest of the cluster by at most two edges, both of them bridges, a
self-avoiding path between two vertices of $B_i$ stays in~$B_i$; so these edges are bridges of $B_i$
separating two of its vertices.
By \cref{lem:switching} the open edges inside $B_i$ are those of~$\omega^{(i)}$. Hence
$c_i-1\le\diambr(B_i)$, the bridge diameter of $B_i$ in~$\omega^{(i)}$. Now $B_i$ is a component of the
$\omega^{(i)}$-open subgraph induced by $C_i\setminus S$, joined to $C_i\cap S$ by the
$\abs{\mathcal K_i}\ge1$ open edges of $\mathcal K_i$. Since $C_i$ is $M$-good, $\diambr(B_i)\le
\abs{\mathcal K_i}M^\sigma$. If $c_i=1$ the claim is $\abs{\mathcal K_i}\ge1$; otherwise $c_i-1\ge
c_i/2$ gives $1/\abs{\mathcal K_i}\le M^\sigma/(c_i-1)\le2M^\sigma/c_i$. Also $c_i-1\le\abs{B_i}-1$,
so $c_i\le M$ and codes indeed take values in $\{\ast\}\cup\{1,\dots,M\}$.

\ref{it:codes-sum}: By~\ref{it:codes-gaps}, the product for a code of a successful run is at most
$\prod_{i\colon c_i\neq\ast}\min\{1,2M^\sigma/c_i\}$. Summing this over all sequences
$c\in(\{\ast\}\cup\{1,\dots,M\})^\ell$ gives $\Theta_M^\ell$. The bound on $\Theta_M$ is the harmonic sum.
\end{proof}

All configurations of a successful run are determined, as far as they matter, by the edges in a fixed
finite ball around~$x$: the cluster $A$ is explored within distance $2n$, the chain moves at most $\nu$
steps from $z\in C_x(\omega'^{(0)})$, each rule reads edges within distance $r_\Phi$ of its vertex, and
every mask lies within distance one of a piece. Conditioning on all other edges, we may therefore treat
the trajectories as elements of a finite product space. On it, the exchanges preserve the total number
of open edges among the $\ell+1$ configurations, so input and output configurations have the same
probability under the product measure. This is the basis of the following density description.

\begin{lemma}[The load]\label{lem:density}
For $x,z\in V$ and a configuration $\omega$ let
\begin{equation*}
  L_x(z,\omega)\coloneqq\frac1n\,\EE\Bigl[\sum_\tau\prod_{i=1}^\ell\Phi_{y_i}(\omega^{(i)},y_{i-1})
  \prod_{j=1}^t\abs{\mathcal K_{i_j}}^{-1}\Bigr]\,,
\end{equation*}
where the expectation is over independent auxiliary configurations $\omega'^{(1)},\dots,\omega'^{(\ell)}$
with law $\PP_{p_c}$, and the sum is over the trajectories $\tau$ of successful runs whose output is
$(\omega,\omega'^{(1)},\dots,\omega'^{(\ell)})$; the factors are evaluated on these trajectories. Then,
for every event~$\mathcal B$,
\begin{equation}\label{eq:density}
  \frac1n\,\PP\bigl(\text{the experiment succeeds and }\omega'^{(0)}\in\mathcal B\bigr)
  =\EE_{p_c}\bigl[L_x(z,\omega)\1_{\mathcal B}(\omega)\bigr]\,.
\end{equation}
Moreover $0\le L_x(z,\omega)\le\Theta_M^\ell/n$; the function $L_x(z,\cdot)$ depends only on the edges in a
finite ball around~$x$; $L_{gx}(gz,g\omega)=L_x(z,\omega)$ for $g\in\Aut(G)$; and $L_x(z,\omega)=0$
unless $z\in C_x(\omega)$ and $\abs{C_x(\omega)}<2n+\nu$.
\end{lemma}

\begin{proof}
A trajectory has probability $\PP_{p_c}(\omega^{(0)})\cdots\PP_{p_c}(\omega^{(\ell)})
\prod_i\Phi_{y_i}(\omega^{(i)},y_{i-1})\prod_j\abs{\mathcal K_{i_j}}^{-1}$ on the finite product space.
The product of the configuration probabilities equals that of the output configurations. Summing over
trajectories with a given output gives~\eqref{eq:density}. For the upper bound, group the trajectories
with a given output by their codes. By \cref{lem:codes}\ref{it:codes-inverse}, the code determines the
input configurations and the seams, so for fixed code the sum over trajectories is a sum over chains,
and $\sum_{\text{chains}}\prod_i\Phi_{y_i}(\omega^{(i)},y_{i-1})\le1$ because these are the
probabilities of the chains given the input configurations. \Cref{lem:codes}\ref{it:codes-sum} bounds
the remaining sum over codes by $\Theta_M^\ell$. Locality and covariance hold because every ingredient is
defined locally and covariantly, and the support property is~\eqref{eq:switch-cluster}.
\end{proof}

\subsection{First and second moment of the load}\label{sec:load}

Let $L(\omega)\coloneqq\sum_zL_x(z,\omega)$ be the \emph{total load} and $\bar L\coloneqq\EE_{p_c}L$; by
invariance these do not depend on~$x$ in law. Write $\pi\coloneqq\pi_n$, and let
\begin{align*}
  T^{(n)}(x,y)&\coloneqq\PP_{p_c}\bigl(y\in C_x,\ n\le\abs{C_x}<2n\bigr)\,,\\
  \chi^{(n)}&\coloneqq\sum_yT^{(n)}(o,y)=\EE_{p_c}\bigl[\abs{C_o};\,n\le\abs{C_o}<2n\bigr]\,.
\end{align*}
The kernel $T^{(n)}$ is symmetric, because $y\in C_x$ implies $C_x=C_y$, and
$n\pi\le\chi^{(n)}\le2n\pi$.

\begin{lemma}[Mean load]\label{lem:mean-load}
The \emph{mean load} $\bar L(x,z)\coloneqq\EE_{p_c}L_x(z,\omega)$ is a kernel with row sums $\bar L$, and
\begin{equation}\label{eq:mean-load}
  \bar L(x,z)=\frac{s^\ell}n\sum_yT^{(n)}(x,y)\,K^\ell(z,y)\,,\qquad
  \pi s^\ell\le\bar L\le2\pi s^\ell\,,\qquad\norm{\bar L(\cdot,\cdot)}\le\bar L\,r^\ell\,.
\end{equation}
\end{lemma}

\begin{proof}
By~\eqref{eq:density}, $\bar L(x,z)$ is $1/n$ times the success probability of the experiment. The $\ell$ steps
of the chain use independent configurations, each succeeding with probability $s$ and moving according
to $K$, so the chain succeeds and ends at $y$ with probability $s^\ell K^\ell(z,y)$; independently, $y\in A$
and $n\le\abs A<2n$ with probability $T^{(n)}(x,y)$. Since $K^\ell$ is doubly stochastic,
$\bar L=\sum_z\bar L(x,z)=s^\ell\chi^{(n)}/n\in[\pi s^\ell,2\pi s^\ell]$. Finally
$\bar L(\cdot,\cdot)=(s^\ell/n)\,T^{(n)}(K^\ell)^\top$, and the Schur test gives
$\norm{T^{(n)}}\le\chi^{(n)}$, so $\norm{\bar L(\cdot,\cdot)}\le(s^\ell/n)\chi^{(n)}r^\ell=\bar Lr^\ell$.
\end{proof}

The mean kernel already has the desired small norm. To turn the load into the output probabilities of a
new map we must divide by a bound on~$L$, and we need that most of the mean is carried by
configurations on which $L$ is not much larger than its mean. This is what the second moment controls.

\begin{lemma}[Second moment of the load]\label{lem:second-moment}
There is a polynomial $Q(\ell,M)\ge1$, with coefficients depending only on~$G$, such that
\begin{equation}\label{eq:second-moment}
  \EE_{p_c}L^2\le Q(\ell,M)\,\pi s^{2\ell}\qquad\text{whenever }n\in\mathcal N,\ \ell M\le n/2\text{ and }
  n\ge(\Theta_M/s)^\ell\,.
\end{equation}
\end{lemma}

\begin{proof}
Expand $L^2=\sum_{z_1,z_2}L_x(z_1,\omega)L_x(z_2,\omega)$, using two independent families of auxiliary
output configurations. Each summand is a weighted sum over pairs $(\tau^1,\tau^2)$ of trajectories of
successful runs with the common distinguished output~$\omega$ and endpoints $z_1,z_2$. We split the pairs
into three classes. For a trajectory with at least one round, let $H^1$ or $H^2$ be its hanging part
(\cref{lem:switching}); it contains the endpoint and has at most $\nu$ vertices. Two hanging parts in the
same cluster are the far sides of bridges, so they are nested or disjoint.

\emph{Class~1: one of the trajectories has no round.} Let $L^\emptyset(\omega)$ be the part of
$L(\omega)$ coming from trajectories without rounds. For them the output equals the input, so
$A=C_x(\omega)$, $n\le\abs A<2n$, $z\in A$, and the chain from $z$ ends in~$A$. Hence
\begin{equation*}
  L^\emptyset(\omega)\le\frac{s^\ell}n\,\1\{n\le\abs{C_x(\omega)}<2n\}\sum_{z\in C_x(\omega)}\
  \sum_{y\in C_x(\omega)}K^\ell(z,y)\le\frac{s^\ell}{n}\,2n=2s^\ell\,,
\end{equation*}
using that the column sums of $K^\ell$ are one. The pairs of class~1 contribute at most
$2\EE_{p_c}[L^\emptyset L]\le4s^\ell\bar L\le8\pi s^{2\ell}$.

\emph{Class~2: the hanging parts are nested or joined by an edge of~$G$.} Fix $\tau^1$, with hanging part
$H^1\ni z_1$. We claim that at most $N_\nu\coloneqq\nu+\nu^2+d\nu^3$ endpoints $z_2$ admit a
trajectory $\tau^2$ in this class. If $H^2\subseteq H^1$, then $z_2\in H^1$: at most $\nu$ choices. If
$H^1\subsetneq H^2$, the first seam of $\tau^2$ is a bridge separating $x$ from $z_1$, and the bridges
between it and $z_1$ lie in $H^2$; so it is among the last $\nu$ such bridges, it determines $H^2$, and
$z_2\in H^2$: at most $\nu^2$ choices. If $H^1$ and $H^2$ are disjoint and joined by an edge, pick a
vertex of $H^2$ adjacent to $H^1$ (at most $d\nu$ choices); the first seam of $\tau^2$ is among the last
$\nu$ bridges separating $x$ from it, and $z_2\in H^2$: at most $d\nu^3$ choices. By
\cref{lem:density}, $L_x(z_2,\omega)\le\Theta_M^\ell/n\le s^\ell$ for each such $z_2$. Hence class~2
contributes at most $\EE_{p_c}[L]N_\nu s^\ell\le2N_\nu\pi s^{2\ell}$.

\emph{Class~3: the hanging parts are disjoint and not joined by any edge of~$G$.} Undo both
trajectories: by \cref{lem:switching} the changes of the distinguished configuration made by $\tau^a$
concern edges with an endpoint in $H^a$, and these two sets of edges are disjoint, so the undoing
operations commute. Let $\omega_{00}$ be the jointly restored distinguished configuration. We claim:
\begin{enumerate}[label=(\alph*)]
  \item $C_x(\omega_{00})=A_{00}\coloneqq C_x(\omega)\setminus(H^1\cup H^2)$, and
  $n-\nu\le\abs{A_{00}}<2n$;
  \item running the switching procedure of $\tau^1$ from $A_{00}$ instead of from
  $A^1\coloneqq C_x(\omega)\setminus H^1$, with the jointly restored configurations and the same chain
  and seams, reproduces the selected indices, pieces, contact sets and masks of~$\tau^1$; the same
  holds for~$\tau^2$. In particular, the first selected cluster of each chain meets $A_{00}$.
\end{enumerate}
For~(a): undoing $\tau^1$ alone restores the input of $\tau^1$, whose cluster of $x$ is $A^1$ by
\cref{lem:switching}; its edges to $H^1$ are closed and remain so, because undoing $\tau^2$ changes only
edges with an endpoint in $H^2$, and no edge joins $H^1$ to~$H^2$. The same holds with the roles
exchanged. Edges with both endpoints in $A_{00}$ are not changed, and $A_{00}$ is connected, because
removing the far sides of two bridges leaves the near side connected. Hence $C_x(\omega_{00})=A_{00}$.
Finally $\abs{A^1}\in[n,2n)$ and $\abs{H^2}\le\nu$.

For~(b) we argue by induction over the rounds of $\tau^1$, writing $S$ for the current set of the
original run and $S_0$ for that of the new run; initially $S_0=S\setminus H^2$, and $H^2\subseteq S$
throughout, since $H^2\subseteq A^1$. Let $i$ be the index selected in the original run and
$e_i=\{a_i,b_i\}$ its seam with $a_i\in S$ and $b_i\in B_i\subseteq H^1$. Since no edge joins $H^1$ to
$H^2$, we have $a_i\notin H^2$, so $C_i\ni a_i$ meets $S_0$, while every larger eligible index misses
$S\supseteq S_0$. Hence the same index is selected. The new piece is the component of $y_i$ in the
$\omega^{(i)}$-open subgraph on $V\setminus S_0=(V\setminus S)\cup H^2$; an open path leaving $B_i$
would have to enter $H^2$ through an edge from $B_i\subseteq H^1$, which does not exist, so the piece is
again~$B_i$. For the same reason the contact sets and the masks of the two runs coincide, and the states
exchanged on the mask agree, because undoing $\tau^2$ changed no edge with an endpoint in~$H^1$.
Afterwards $S_0\cup B_i=(S\cup B_i)\setminus H^2$. The stopping rules agree because $z_1\in H^1$ is not
in~$H^2$. The first selected cluster meets $S_0=A_{00}$ by the argument just given.

By~(b), a pair $(\tau^1,\tau^2)$ of class~3 is determined by the jointly restored configurations
together with the two chains and the two seam sequences: rerunning both procedures from $A_{00}$ and
applying their disjoint exchanges recovers the output. The exchanges preserve the total number of open
edges, so the weight of the pair equals the probability of the restored configurations, times the
chain probabilities and the factors $\prod\abs{\mathcal K}^{-1}$ of the two procedures run from
$A_{00}$. For fixed restored configurations and chains, the selected indices, the pieces and the contact
sets of each procedure do not depend on the seam choices, because every round adds the whole piece to
the current set; only the masks and the output depend on the seams. Hence the factors
$\prod\abs{\mathcal K}^{-1}$ are the probabilities of the uniform seam choices of each run, their sum
over the seam sequences of one run is one, and the sum over pairs of seam sequences is at most the
product of these two sums. Dropping
all conditions except~(a) and the fact that each chain succeeds and meets $A_{00}$, we obtain that
class~3 contributes at most
\begin{equation*}
  \frac1{n^2}\sum_{z_1,z_2}\EE_{p_c}\bigl[\1\{n-\nu\le\abs{C_x}<2n\}\,
  \mathrm{hit}_{z_1}(C_x)\,\mathrm{hit}_{z_2}(C_x)\bigr]\,,
\end{equation*}
where, for a fixed set $D$, $\mathrm{hit}_z(D)$ is the probability that a chain of $\ell$ independent
steps of $\Phi$ from $z$ succeeds and that one of its clusters $C_i$ meets~$D$. To bound it, let every
vertex $z$ send unit mass to each vertex of $C_1\cup\dots\cup C_\ell$ when its chain succeeds. A successful
chain has probability $s^\ell$ and its clusters have at most $\nu$ vertices in total, so each vertex sends
expected mass at most $\nu s^\ell$, and by the mass-transport principle each vertex receives at most
$\nu s^\ell$. A union bound over the vertices of $D$ gives
\begin{equation*}
  \sum_z\mathrm{hit}_z(D)\le\nu s^\ell\abs D\,.
\end{equation*}
Since $\nu\le n/2$, class~3 therefore contributes at most
\begin{equation*}
  \frac{\nu^2s^{2\ell}}{n^2}\,\EE_{p_c}\bigl[\abs{C_x}^2;\,n/2\le\abs{C_x}<2n\bigr]
  \le4\nu^2s^{2\ell}\,\PP_{p_c}(n/2\le\abs{C_o}<2n)\le4(1+2^\beta)\nu^2\pi s^{2\ell}\,,
\end{equation*}
by~\eqref{eq:scales}. Adding the three classes gives~\eqref{eq:second-moment} with
$Q(\ell,M)\coloneqq8+2N_\nu+10\nu^2$, since $4(1+2^\beta)\le10$.
\end{proof}

\subsection{One amplification step and the iteration}\label{sec:iteration}

By \cref{prop:good} there is $M_\varepsilon$ such that the bound of \cref{prop:good} holds for all
$M\ge M_\varepsilon$.

\begin{lemma}[One amplification step]\label{lem:one-step}
Let $\Phi$ be a local endpoint map with cap $M\ge2$, success probability $s$ and kernel norm~$r$. Let
$\ell\in\N$ and $n\in\mathcal N$ satisfy
\begin{equation}\label{eq:one-step-conditions}
  \ell M\le n/2\,,\qquad n\ge(\Theta_M/s)^\ell\,,\qquad 2n\ge M_\varepsilon\,,\qquad
  16\,Q(\ell,M)\,n^\beta\e^{-n^\varepsilon}\le1\,.
\end{equation}
Then there is a local endpoint map $\Phi'$ with cap $M'\coloneqq2n+\ell M$, success probability
$s'\ge\pi_n/(8Q(\ell,M))$ and kernel norm $r'\le2r^\ell$.
\end{lemma}

\begin{proof}
Put $b\coloneqq4Q(\ell,M)s^\ell$. By \cref{lem:mean-load,lem:second-moment},
$\EE_{p_c}[L;L>b]\le\EE_{p_c}L^2/b\le\pi s^\ell/4\le\bar L/4$. The load vanishes unless
$\abs{C_x}<2n+\nu\le M'$, so by \cref{prop:good}, $M'\ge M_\varepsilon$ and $\pi_n\ge n^{-\beta}$,
\begin{equation*}
  \EE_{p_c}\bigl[L;\,L\le b,\ C_x\text{ not $M'$-good}\bigr]\le b\,\e^{-(M')^\varepsilon}
  \le\bar L\cdot\frac{4Q(\ell,M)}{\pi_n}\,\e^{-n^\varepsilon}\le\frac{\bar L}4\,,
\end{equation*}
by the last condition in~\eqref{eq:one-step-conditions}. Let $\mathcal R$ be the event that $L\le b$
and $C_x$ is $M'$-good, and define
\begin{equation*}
  \Phi'_x(\omega,z)\coloneqq\1_{\mathcal R}(\omega)\,\frac{L_x(z,\omega)}{b}\,.
\end{equation*}
By \cref{lem:density} this is a local endpoint map with cap~$M'$: on $\mathcal R$ the outputs have total
probability $L/b\le1$, every output lies in $C_x$, and $\abs{C_x}<M'$ with $C_x$ being $M'$-good. Its
unconditional endpoint kernel $F'_{p_c}(x,z)=\EE_{p_c}[L_x(z,\omega);\mathcal R]/b$ satisfies
$0\le bF'_{p_c}(x,z)\le\bar L(x,z)$, and its row sums are $s'=\EE_{p_c}[L;\mathcal R]/b\ge\bar L/(2b)\ge\pi_n/(8Q(\ell,M))$. Hence
$r'=\norm{F'_{p_c}}/s'\le\norm{\bar L(\cdot,\cdot)}/(\bar L/2)\le2r^\ell$ by \cref{lem:kernels}\ref{it:domination}
and~\eqref{eq:mean-load}.
\end{proof}

The new map $\Phi'$ consists of fixed functions $\Phi'_x(\omega,\cdot)$ of the configuration. Like every
local endpoint map, it can be evaluated on configurations of any parameter; this is how the maps are used
in \cref{sec:moments}.

\begin{proof}[Proof of \cref{prop:amplification}]
We measure a map $\Phi$ with cap $M$ and success probability $s$ by
\begin{equation*}
  \operatorname{cost}(\Phi)\coloneqq\log(1/s)+(\sigma+\xi)\log M\,,
\end{equation*}
where $\xi>0$ is fixed so small that, by~\eqref{eq:beta-sigma},
\begin{equation*}
  c_\xi\coloneqq\frac{\beta+\xi+(\sigma+\xi)(1+\xi)}{1-\xi}<1\,.
\end{equation*}
Fix the constant $C$ in $\Theta_M\le CM^\sigma\log M$, and choose $M_1\ge M_\varepsilon$ so large that
$\log(2C\log M)\le\xi\log M$ for $M\ge M_1$. By \cref{lem:start} there is a map $\Phi_0$ with cap
$M\ge M_1$ and kernel norm $r<\e^{-1}$.

\emph{One step.} Let $\Phi$ have cap $M\ge M_1$, success probability $s$ and kernel norm
$r<\e^{-1}$, and put $a\coloneqq\lceil CM^\sigma\log M/s\rceil$. Then $a\ge\max\{2,\Theta_M/s\}$ and
$\log a\le\log(1/s)+\sigma\log M+\log(2C\log M)\le\operatorname{cost}(\Phi)$. For $n\in\mathcal N$ put
$\ell\coloneqq\lfloor\log n/\log a\rfloor$, so that $a^\ell\le n<a^{\ell+1}$. As $n\to\infty$ through
$\mathcal N$, we have $\ell\to\infty$ and $\ell\le\log n/\log2$; hence all conditions
in~\eqref{eq:one-step-conditions} hold for large~$n$, using that $Q(\ell,M)$ is polynomial in~$\ell$ and
$n\ge a^\ell\ge(\Theta_M/s)^\ell$. \Cref{lem:one-step} gives a map $\Phi'$ with cap $M'$, success
probability $s'$ and kernel norm $r'$. By $\pi_n\ge n^{-\beta}$ and
$\log n<(\ell+1)\operatorname{cost}(\Phi)$,
\begin{align*}
  \log(1/s')&\le\beta\log n+\log(8Q(\ell,M))\le\beta(\ell+1)\operatorname{cost}(\Phi)+O_M(\log(\ell+1))\,,\\
  \log M'&=\log(2n+\ell M)\le\log n+1\le(\ell+1)\operatorname{cost}(\Phi)+1\,,\\
  \log(1/r')&\ge\ell\log(1/r)-\log2\,.
\end{align*}
Since $\operatorname{cost}(\Phi)>0$ and $\log(1/r)>1$ are fixed, for $n$ large enough this gives
$\log(1/s')\le\ell(\beta+\xi)\operatorname{cost}(\Phi)$, $\log M'\le\ell(1+\xi)\operatorname{cost}(\Phi)$
and $\log(1/r')\ge\ell(1-\xi)\log(1/r)\ge2\log(1/r)$, and therefore
\begin{equation*}
  \frac{\operatorname{cost}(\Phi')}{\log(1/r')}
  \le\frac{\ell\operatorname{cost}(\Phi)\,\bigl(\beta+\xi+(\sigma+\xi)(1+\xi)\bigr)}{\ell(1-\xi)\log(1/r)}
  =c_\xi\,\frac{\operatorname{cost}(\Phi)}{\log(1/r)}\,.
\end{equation*}
Moreover $\Phi'$ has cap $M'\ge2n\ge M_1$ and $r'<\e^{-2}$.

\emph{Iteration.} Starting from $\Phi_0$ and applying the step repeatedly, we obtain maps
$\Phi_0,\Phi_1,\dots$ with kernel norms $r_j<\e^{-1}$ and
\begin{equation*}
  \frac{\operatorname{cost}(\Phi_j)}{\log(1/r_j)}\le c_\xi^j\,\frac{\operatorname{cost}(\Phi_0)}{\log(1/r_0)}
  \longrightarrow0\,.
\end{equation*}
Since $\log(1/s)+\log M\le\max\{1,(\sigma+\xi)^{-1}\}\operatorname{cost}(\Phi)$ for every map, the
efficiency ratio $\operatorname{eff}(\Phi_j)$ tends to zero.
\end{proof}
